\uuid{szfL}
\titre{Changements de variables}
\auteur{hanani ; mzari}
\organisation{exo7}
\licence{CC BY-NC-SA 4.0}
\datecreate{2020-01-20}
\niveau{L1}
\module{Analyse}
\chapitre{Calcul d'intégrales}
\sousChapitre{Autre}
\difficulte{3}
\theme{Calculs d'intégrales, changements de variables}
\exo7qcmid{lille-2.integrales.018}

\contenu{
\begin{question}
\qclassification{127.99}
\qtags{motcle=changements de variables}
On pose $\displaystyle I=\int _{\pi/6}^{\pi/3}\frac{\cos x}{\sin x}\mathrm{d}x$ et $\displaystyle J=\int _{\pi/6}^{\pi/3}\frac{\sin x}{\cos x}\mathrm{d}x$. Parmi les affirmations suivantes, cocher celles qui sont vraies :
\begin{answers}
\good{$\displaystyle I=J$.}
\bad{$\displaystyle I=\frac{1}{J}$.}
\good{$\displaystyle I=\frac{1}{2}\ln 3$.}
\bad{$\displaystyle J=-\ln \sqrt{3}$.}
\end{answers}
\begin{explanations}
Avec $t=\pi/2- x\Rightarrow \mathrm{d}t=-\mathrm{d}x$, on obtient : $\displaystyle I=-\int _{\pi/3}^{\pi/6}\frac{\sin t}{\cos t}\mathrm{d}t=J$.
\vskip0mm
Avec $t=\sin x\Rightarrow \mathrm{d}t=\cos x\, \mathrm{d}x$, on obtient : $\displaystyle I=\int _{1/2}^{\sqrt{3}/2}\frac{\mathrm{d}t}{t}=\ln \sqrt{3}$.
\end{explanations}
\end{question}
}
